\subsection{One finite-ramification hypothesis for the Legendre system}

\begin{claim}{D.30}{Proved}\label{clm:D30}
Let
\[
 U=\mathbf P^1_{\mathbf Q}\setminus\{0,1,\infty\},\qquad
 \mathcal E/U:\ y^2=x(x-1)(x-t),
\]
and let $s:G_{\mathbf Q}\to\pi_1(U)$ be a continuous
noncuspidal section. Write $\rho_{s,\ell}$ for the
two-dimensional rational $\ell$-adic representation obtained
by pulling back the Legendre Tate module along $s$.
Then every $\rho_{s,\ell}$ is absolutely irreducible,
without a finite-ramification assumption.

The following conditions are equivalent:
\begin{enumerate}
\item For at least one prime $\ell\ge5$, $\rho_{s,\ell}$
is ramified at only finitely many rational primes.
\item The associated local sections are ordinary integral
point sections at all but finitely many finite places:
$x_v\in U(\mathbf Z_v)$ there.
\item There is one finite set $S$ such that, for every
coefficient prime $\ell$, $\rho_{s,\ell}$ is unramified
at every $v\notin S\cup\{\ell\}$.
\end{enumerate}
Under these conditions the rational representation system
is the Tate-module system of an elliptic curve over
$\mathbf Q$. In particular its representations with
coefficient prime at least five are modular of weight two
and trivial nebentype.

The assertion does not prove condition \textup{(1)} for
an arbitrary section, and it does not identify $s$ with
a rational point section. Condition \textup{(2)} permits
finitely many cuspidal local sections. D.5 and Stix's
Corollary~6 require ordinary local points at every finite
place, so they cannot be applied without a further argument
eliminating those remaining cuspidal places.
\end{claim}

\begin{proof}
The Weil pairing gives
$\det(\rho_{s,\ell})=\chi_\ell$. D.6 supplies the ordinary
point or cuspidal classification at each finite place;
D.25 supplies a density-one set of ordinary integral places.
All appearances of Hodge--Tate weights below use the
convention that $\mathbf Q_\ell(1)$ has weight $-1$.

\emph{The local representation at its coefficient prime.}
For an ordinary local section at $\ell$, the representation
is the Tate module of the actual elliptic curve
$\mathcal E_{x_\ell}/\mathbf Q_\ell$. It is potentially
semistable with weights $0,-1$.
For a cuspidal local section, the Legendre family has
potentially multiplicative reduction. After a quadratic
twist, its cusp-decomposition representation is an extension
\[
 0\longrightarrow\mathbf Q_\ell(1)\longrightarrow V
   \longrightarrow\mathbf Q_\ell\longrightarrow0.
\]
Pullback along any continuous cuspidal splitting gives a
class in $H^1(\mathbf Q_\ell,\mathbf Q_\ell(1))$.
Every such extension is semistable, including the split
one and classes not supplied by a chosen rational tangent
vector: \citet[II.4.4, pp.~18--19]{D_Berger2002} proves
this from Kummer theory and the Tate-curve periods.
Undoing the quadratic twist preserves potential
semistability and the two weights. The finite twist at
the infinity cusp is retained; it is not silently discarded
as an unramified character.

\emph{A character lemma, with arbitrary finite coefficient field.}
Let $E/\mathbf Q_\ell$ be finite, and suppose that
$\psi:G_{\mathbf Q}\to E^\times$ is continuous and
Hodge--Tate at $\ell$. We prove
\[
                 \psi=\eta\chi_\ell^a,\qquad
 a\in\mathbf Z,\quad\eta\text{ of finite order}.
 \tag{D.30.1}
\]
Compactness puts its image in $\mathcal O_E^\times$.
Choose $h$ such that the $h$-th powers of this unit group
lie in a torsion-free pro-$\ell$ group of principal units.
Kronecker--Weber identifies
\[
 G_{\mathbf Q}^{\mathrm{ab}}
  \simeq\widehat{\mathbf Z}^{\,\times}
  =\prod_q\mathbf Z_q^\times.
\]
For $q\ne\ell$, a continuous map from $\mathbf Z_q^\times$
to a torsion-free pro-$\ell$ group kills its finite factor
and its open pro-$q$ subgroup. Thus $\psi^h$ kills every
such factor. Their direct sum is dense in their product,
so $\psi^h$ factors through $\mathbf Z_\ell^\times$.

Apply Tate's local algebraicity theorem
\citep[III, Appendix A.7, Theorem~3, pp.~III-52--53]{D_Serre1989}
to the underlying $\mathbf Q_\ell$-representation of the
$E$-line. Its inertia action is semisimple, since it is
given by multiplication by elements of a characteristic-zero
field. The resulting algebraic inertia action is a morphism
\[
 \mathbf G_m/\mathbf Q_\ell
   \longrightarrow\operatorname{GL}_{\mathbf Q_\ell}(E).
\]
It centralizes $E$, as it does so on the Zariski-dense
open subgroup of local units. It therefore factors through
$\operatorname{Res}_{E/\mathbf Q_\ell}\mathbf G_m$.
A $\mathbf Q_\ell$-defined cocharacter of that torus has
the same integer exponent at every coefficient embedding,
because those embeddings form one Galois orbit. Hence,
with the cyclotomic convention fixing the sign,
$\psi\chi_\ell^{-a}$ is trivial on an open subgroup of
$I_\ell$ for one integer $a$.

For $\eta=\psi\chi_\ell^{-a}$, the preceding power argument
shows that $\eta^h$ factors through $\mathbf Z_\ell^\times$
and is trivial on an open subgroup there, so it has finite
image. The $h$-th power map on $E^\times$ has finite
fibres; consequently $\eta$ itself has finite image.
This proves \textup{(D.30.1)} without assuming finite
ramification of $\psi$ in advance.
This coefficient-field argument is necessary: an invariant
line of $\rho_{s,\ell}$ need not initially be
$\mathbf Q_\ell$-rational.

\emph{Absolute irreducibility.}
If $\rho_{s,\ell}$ were reducible over
$\overline{\mathbf Q}_\ell$, an invariant line would be
defined over a finite extension $E$.
Its character and the quotient character are Hodge--Tate;
their weights add to $0,-1$
\citep[Proposition~2.4.4]{D_BrinonConrad2009}.
By \textup{(D.30.1)} and the determinant, their
semisimplification, after interchanging them, is
\[
                    \eta\oplus\eta^{-1}\chi_\ell
\]
for a finite-order character $\eta$.
The primes splitting in the finite abelian extension
killing $\eta$ have positive density. Their intersection
with the density-one ordinary integral set of D.25
still has positive density. At each such prime
$v\ne\ell$, the semisimplification has trace $1+v$.
The same trace belongs to a good-reduction elliptic curve,
and must satisfy $|a_v|\le2\sqrt v$. Since
$v+1>2\sqrt v$, this is impossible.
Only the one-dimensional characters were shown finitely
ramified in this argument, not the original representation.

\emph{The precise modularity theorem.}
Assume \textup{(1)}, and fix the indicated $\ell\ge5$.
The representation is absolutely irreducible, is odd by
its cyclotomic determinant, and is potentially semistable
at $\ell$ with distinct weights. These are exactly the
hypotheses of \citet[Theorem~1.0.4]{D_Pan2022}, including
the finite-ramification hypothesis incorporated from
Conjecture~1.0.1. No further residual condition is required
when $\ell\ge5$; the exceptional condition in that statement
concerns $\ell=3$.

Thus $\rho_{s,\ell}$ comes from a cuspidal eigenform up
to twist. If the twisting character is denoted $\psi$,
it is a de Rham subobject of
$\operatorname{Hom}(V_f,\rho_{s,\ell})$.
The character lemma makes it finite order times a Tate
twist. The two consecutive weights force classical weight
two and zero Tate twist; absorb the finite-order twist
into the eigenform. Its nebentype is trivial because
the determinant is $\chi_\ell$.
Outside the level and the coefficient prime, the Frobenius
roots are Weil numbers of weight one, so their trace has
absolute value at most $2\sqrt v$.
This purity is an output of modularity with its finite
ramification hypothesis, not a replacement for that
hypothesis.

\emph{Cofinite integrality.}
Enlarge the finite exceptional set to contain two, $\ell$,
the ramified primes, and the modular level. Let $v$ lie
outside it.
If $s_v$ is ordinary, the actual elliptic curve
$\mathcal E_{x_v}/\mathbf Q_v$ has unramified Tate module.
N\'eron--Ogg--Shafarevich gives good reduction
\citep[IV.1.3]{D_Serre1989}.
But
\[
 j(t)=256\,\frac{(t^2-t+1)^3}{t^2(1-t)^2}.
 \tag{D.30.2}
\]
At an odd place, if $v(t)>0$ or $v(1-t)>0$, the numerator
is a unit and $v(j(t))<0$. If $v(t)<0$, the valuation
is $2v(t)<0$. Thus every value outside $U(\mathbf Z_v)$
is not even potentially good, contradicting good reduction.

If $s_v$ is cuspidal, its representation has the form
\[
 \delta_v\otimes
 \begin{pmatrix}\chi_\ell& *\\0&1\end{pmatrix}
\]
for a quadratic character $\delta_v$. This character
need not be unramified a priori. However the whole
representation and $\chi_\ell$ are unramified at the
chosen $v$, so the diagonal entries force $\delta_v$
to be unramified. The trace is then
$\delta_v(\operatorname{Frob}_v)(v+1)=\pm(v+1)$,
contrary to modular weight-one purity.
This excludes the locally split cuspidal case as well:
no nonzero unipotent inertia operator is assumed.
Hence \textup{(1)} implies \textup{(2)}.

Conversely, at an odd place with $x_v\in U(\mathbf Z_v)$,
the Legendre equation defines an elliptic scheme over
$\mathbf Z_v$. Its Tate module is unramified for every
coefficient prime different from $v$.
Enlarging the exceptional set in \textup{(2)} by two
therefore gives \textup{(3)}, which immediately implies
\textup{(1)}. This proves the equivalence.
The same cofinite-integrality conclusion transfers to
a fixed finite \'etale cover of $U$ by the finite
normalization argument of D.25--D.26, after removing
finitely many model primes.

\emph{One elliptic realization of the rational system.}
Let $f$ be the normalized newform just obtained.
Its good-prime coefficients on the density-one integral
set are the rational integer traces of
$\mathcal E_{x_v}$.
For any embedding $\tau$ of its Hecke field, the
representations of $f$ and $f^\tau$, embedded into one
algebraic closure of $\mathbf Q_\ell$, have equal traces
on that set. Such Frobenius elements are dense in every
finite Galois quotient by Chebotarev. Continuity and
semisimplicity give equal characters, hence equal
good-prime eigenvalues everywhere.
Newform multiplicity one away from the level
\citep[Theorem~9.4 and the paragraph after Example~9.5]{D_Stein2007}
gives $f^\tau=f$. Here the usual Hecke recurrences
recover coprime composite indices from the good primes
and the nebentype. Thus the Hecke field is $\mathbf Q$.

The Eichler--Shimura construction gives an elliptic curve
$E/\mathbf Q$ for this rational weight-two newform
\citep[Proposition~11.21]{D_MilneMF2017};
only existence up to isogeny is used.
For every other coefficient prime $r$, the representations
$V_r(E)$ and $\rho_{s,r}$ have the same rational integer
traces on the common density-one set of integral good
primes. Absolute irreducibility was already proved for
$\rho_{s,r}$. Chebotarev and Brauer--Nesbitt yield
\[
                       \rho_{s,r}\simeq V_r(E).
\]
This proves the rational-system assertion.
It does not recover the missing nonabelian information
in the section, and equality of local Tate representations
has not been used to identify the parameters $x_v$.
\end{proof}

\proofparagraph{A genuine cuspidal warning about strict compatibility.}
Even a continuous global section can have different
Weil--Deligne monodromy ranks at two coefficient primes
if it is cuspidal.
Let $e_5\in\widehat{\mathbf Z}$ have component one in
$\mathbf Z_5$ and zero at all other primes.
Twist the tangent-one splitting at $t=0$ by the continuous
Kummer cocycle $e_5\kappa(13)$.
This is a continuous global splitting into the rational
cusp decomposition group and hence a section of $\pi_1(U)$.

At 13 the nodal splitting character is trivial:
$5^2=-1$ modulo 13 and Hensel's lemma supplies the square
root in $\mathbf Q_{13}$. From \textup{(D.30.2)} and
$q=j^{-1}+O(j^{-2})$, the formal Tate parameter is
\[
              q(t)=t^2u(t),\qquad u(0)=1/256.
\]
The reference tangent section thus has an unramified
extension class at 13 for coefficient primes different
from 13. The twist changes the coefficient-$r$ class by
$2(e_5)_r\kappa_r(13)$.
Its monodromy at 13 is consequently nonzero for $r=5$
and zero for $r=7$, whereas both semisimplified
Frobenius polynomials are
\[
                         (T-1)(T-13).
\]
Zero versus nonzero monodromy is an isomorphism invariant.
Thus matching semisimple Frobenius polynomials does not
automatically give strict compatibility at ramified places.
The exact certificate
\path{code/D_legendre_modularity_checks.py} and its JSON
check the leading term, the Hensel residue, six compatible
CRT levels, the tame matrix relation and the two ranks.
It is a certificate for this cuspidal example, not a
noncuspidal counterexample or a finite-ramification test.

\proofparagraph{The remaining finite-ramification assertion.}
General semisimple representations have density-zero
ramification, not necessarily finite ramification
\citep[Theorem~1]{D_KhareRajan2000}.
The inspected primary preprint
\citep[Section~2.2, Main Theorem; Section~2.4, Corollary~14]
{D_KhareLarsenRamakrishna2003}
constructs infinitely ramified lifts with potentially
semistable coefficient-prime behavior and integral pure
good-prime polynomials, and density-one compatible pairs.
Its full proof permits weight one: the newly prescribed
primes are chosen after the bound
$v>(\ell^m/2)^{2/k}$, which applies to $k=1$.
It expressly distinguishes density-one compatibility from
matching Weil--Deligne parameters at ramified primes.
The exact numbering and weight range refer to the opened
2003 author preprint; the final journal text was not
retrieved and is not silently substituted for that source.

Those examples are not asserted to arise from tripod
sections. For the present route the missing assertion is
still that a noncuspidal section has one finitely ramified
$\rho_{s,\ell}$ with $\ell\ge5$.
Neither D.25's density statement nor the local
point/cusp classification establishes it.
Moreover, even the equivalent cofinite-integrality
conclusion leaves the finitely many cuspidal local places
and the full nonabelian section comparison unresolved.
The independent restrictions in D.27--D.29 are unchanged.
